
\begin{abstract}
While reinforcement learning algorithms have made significant progress in solving multi-armed bandit problems, they often lack biological plausibility in architecture and dynamics.
Here, we propose a bio-inspired neural model based on interacting populations of rate neurons, drawing inspiration from the orbitofrontal cortex and anterior cingulate cortex.
Our model reports robust performance across various stochastic bandit problems, matching the effectiveness of standard algorithms such as Thompson Sampling and UCB.
Notably, the model exhibits adaptive behavior: employing greedy strategies in low-uncertainty situations while increasing exploratory behavior as uncertainty rises.
Through evolutionary optimization, the model's hyperparameters converged to values that align with the principles of synaptic mechanisms, particularly in terms of synapse-dependent neural activity and learning rate adaptation.
These findings suggest that biologically-inspired computational architectures can achieve competitive performance while providing insights into neural mechanisms of decision-making under uncertainty.
\end{abstract}
